% !TEX root = ../main.tex
\part{Organização de Esquadrão}

Um esquadrão bem montado é a diferença entre uma força equilibrada e um bando de soldados condenados. Esta parte explica como construir, equipar e especializar suas unidades de infantaria --- as escolhas que você faz antes da partida começar são tão importantes quanto qualquer decisão tomada durante o combate.

\section{Slots de Equipamento}

Para agilizar a preparação da partida, Vanguarda 44 não utiliza contagem complexa de pontos por miniatura dentro do esquadrão. Em vez disso, utiliza-se o sistema de Slots de Dado: um limite fixo de 10 ``espaços de dado'' por Esquadrão de Infantaria.

Esse sistema de Slots vale apenas para a composição interna de cada Esquadrão de Infantaria. No nível do Pelotão --- quantos esquadrões, times de suporte e veículos você inclui na sua força --- Vanguarda 44 usa um sistema de pontos simples, explicado na Parte VII, Construção de Pelotão. As duas camadas não se misturam: dentro do esquadrão, nunca há contagem de pontos; entre unidades do Pelotão, sim.

\subsection{Custo de Dados por Arma}

\begin{center}
\begin{tabular}{lc}
\toprule
Arma & Custo em Slots \\
\midrule
Rifle & 1 \\
Pistola (Médico/Oficial) & 1 \\
Submetralhadora (SMG) & 2 \\
Metralhadora Leve (LMG / BAR) & 4 (ou 3 se for BAR americano) \\
\bottomrule
\end{tabular}
\end{center}

\subsection{A Estrutura Base (Padrão)}

Todo Esquadrão de Infantaria começa com esta formação equilibrada, ocupando exatamente os 10 espaços de dado disponíveis:

\begin{center}
\begin{tabular}{p{4.4cm}p{4.8cm}p{2.4cm}}
\toprule
Função & Armamento & Custo (slots) \\
\midrule
1 Sargento (Líder) & Submetralhadora (SMG) & 2 \\
1 Equipe de Metralhadora Leve & 1 Atirador (LMG) + 1 Municiador (Rifle) & 4 + 1 = 5 \\
3 Fuzileiros & Rifle & 3 \\
\bottomrule
\end{tabular}
\end{center}

\textbf{Total}: 6 homens, 10 espaços de dado.

\begin{nota}
O limite de 10 espaços de dado é uma regra rígida: o esquadrão sempre deve ter exatamente o tamanho que caiba nesses 10 espaços, nunca mais soldados do que isso permite. O número de homens em campo varia conforme o equipamento escolhido --- trocar Fuzileiros por armas mais caras reduz o efetivo, mas o total de espaços de dado ocupados é sempre 10.
\end{nota}

\subsection{Coesão de Unidade}

Ao mover um Esquadrão, nenhuma miniatura pode terminar o movimento a mais de 2,5cm (1 polegada) de outra miniatura do mesmo esquadrão. Esta é a Coesão de Unidade padrão, referenciada em outras partes do livro (Fogo e Movimento, Comando de Pelotão, Ocupação Parcial de Edifícios).

\subsection{O Dilema Tático (Qualidade vs. Quantidade)}

Como o limite é fixo em 10 dados, o jogador escolhe entre ter mais soldados (mais ``vidas'') ou armas melhores (mais dano, mas mais frágeis).

\begin{center}
\begin{tabular}{p{4.4cm}p{2.6cm}p{5.4cm}}
\toprule
Composição & Dados / Vidas & Vantagem \\
\midrule
Pelotão Padrão: 1 Sargento SMG + 8 Rifles & 10 / 9 & Bom equilíbrio para aguentar baixas e capturar objetivos \\
Parede de Carne: 10 Rifles & 10 / 10 & Máxima durabilidade, difícil de remover do objetivo \\
Esquadrão de Choque: 5 SMG & 10 / 5 & Devastador de perto e móvel; muito frágil (canhão de vidro) \\
Ninho de Metralhadora: LMG + Sargento SMG + 4 Rifles & 10 / 6 & Poder de fogo de longo alcance e supressão \\
\bottomrule
\end{tabular}
\end{center}

\subsection{Trocas Táticas (Especialistas)}

Antes da partida, o jogador pode dispensar Fuzileiros para incluir armas ou especialistas. Cada Squad pode ter até 1 Especialista (Médico, Engenheiro) que não conta para o limite de 10 dados, mas ocupa 1 vaga de Fuzileiro.

\begin{center}
\begin{tabular}{p{4.2cm}p{3.4cm}p{4.8cm}}
\toprule
Adicionar... & Custo (sai do Squad) & Função Tática \\
\midrule
1 Equipe Bazooka/Panzerschreck (2 homens) & Remove 2 Fuzileiros & Anti-Tanque: perde volume de tiro contra infantaria, ganha capacidade de destruir blindados \\
1 Médico de Campo (1 homem) & Remove 1 Fuzileiro (Especialista) & Sobrevivência: chance de salvar uma baixa por turno \\
2 Soldados de Assalto (SMG) & Remove 2 Fuzileiros & Combate Urbano: troca alcance por dados extras de SMG \\
Até 2 Engenheiros & Substituem Fuzileiros & Cortar arame, desarmar minas, carga de demolição \\
\bottomrule
\end{tabular}
\end{center}

\section{Perfis de Armas --- SMG vs. LMG/BAR}

A diferença entre a Submetralhadora e a Metralhadora Leve não é só o número de dados --- é o estilo de jogo que cada uma favorece. A SMG é a arma do Assalto; a LMG/BAR é a arma de Apoio.

\begin{center}
\begin{tabular}{p{3.6cm}p{4.4cm}p{4.4cm}}
\toprule
Característica & SMG (Thompson/MP40) & LMG / BAR \\
\midrule
Dados & 2 & 3 (BAR) ou 4 (LMG completa) \\
Estilo de jogo & Agressivo: correr e limpar casas & Defensivo: ficar em cobertura, negar área \\
Penalidade de Movimento & Nenhuma (sempre 4+) & -1 se a unidade se mover (vira 5+) \\
Alcance & Curto (20cm) & Longo (60--80cm) \\
Stress causado & Sobe 1 nível & Sobe 2 níveis (regra de Supressão) \\
\bottomrule
\end{tabular}
\end{center}

\begin{nota}
Se a LMG tentar agir como uma SMG (correndo para frente), ela perde grande parte de sua eficácia. Da mesma forma, a SMG não compensa em tiroteios de longa distância.
\end{nota}

\section{O Médico de Combate}

O Médico é um especialista que foca na sobrevivência e no bem-estar psicológico da tropa. Ele ocupa uma vaga de Especialista no Esquadrão (contribuindo com 1 dado de Pistola a curto alcance) ou atua como Unidade de Suporte independente junto ao Tenente.

\textbf{Armamento}: Pistola (Alcance 10cm, 1 Dado). Luta apenas em defesa pessoal.

\textbf{Traço ``Não Combatente''}: o inimigo não pode alvejar o Médico intencionalmente se houver outra miniatura do mesmo esquadrão mais próxima.

\subsection{Presença Calmante (Bônus Passivo)}

No início da ativação do Esquadrão em que o Médico está integrado (antes de mover ou atirar), a unidade reduz automaticamente em 1 o seu Nível de Stress. Este efeito é automático e não gasta PC, mas o Médico deve estar vivo.

\subsection{Primeiros Socorros (Save Roll)}

O Médico pode tentar estabilizar ferimentos fatais em qualquer aliado a até 15cm de distância --- não precisa ser do mesmo esquadrão.

\textbf{O Teste}: imediatamente após o inimigo confirmar uma baixa no Passo B (Jogada de Dano), role 1D6.

\textbf{Resultado 5 ou 6}: o soldado é estabilizado. A miniatura permanece em jogo no estado ``Deitado/Down'' --- não age no próximo turno, mas volta no seguinte.

\textbf{Resultado 1 a 4}: o ferimento é grave demais. O soldado é removido normalmente.

\textbf{Limitação}: o Médico só pode realizar 1 Teste de Salvamento por turno inimigo.

\subsection{Limitações (Dano Catastrófico)}

A jogada de Salvar NÃO é permitida se a baixa foi causada por:

\begin{itemize}[leftmargin=1.2cm]
  \item Armas Pesadas (HE): explosões de artilharia, morteiros ou tanques.
  \item Snipers usando a regra de ``Tiro Selecionado'' (morte instantânea).
  \item Combate Corpo a Corpo (Assalto).
  \item Lança-chamas.
\end{itemize}

\section{Engenheiros de Combate (Pioneiros)}

Engenheiros são especialistas que lidam com obstáculos estáticos. Podem ser incluídos em Esquadrões (máximo 2), substituindo Fuzileiros. Equipamento: Submetralhadora (SMG --- 2 dados) + Kit de Engenharia.

\subsection{Cortar Arame (Breaching)}

Se um Engenheiro estiver em contato de base com uma seção de arame farpado, ele pode gastar sua ação de disparo para remover aquela seção do jogo, sem necessidade de teste.

\subsection{Desarmar Minas}

O Engenheiro deve mover-se para contato com o marcador de mina (sem entrar nele).

O Sargento dá a ordem de ``Engenharia'' (gasta 1 PC).

Role 1D6: 4, 5 ou 6 --- mina neutralizada. 2 ou 3 --- falha técnica, tente de novo no próximo turno. 1 (Falha Crítica) --- o Engenheiro é removido como baixa e a mina explode.

\subsection{Carga de Demolição (Satchel Charge)}

Arma de Uso Único. Alcance de arremesso de 10cm ou contato de base. Contra Edifícios/Bunkers: se acertar (4+), causa 2D6 Acertos Automáticos em quem estiver dentro, ignora a Cobertura Pesada do prédio, e sobe o Nível de Stress em 3 automaticamente.
